Celery es un sistema distribuido de cola de tareas para Python,
ampliamente usado en aplicaciones Django para descargar al servidor web
del trabajo costoso (procesamiento de imágenes, \acs{ocr}, envío de
correo) y permitir su ejecución asíncrona en procesos \textit{worker}
separados \cite{celery}. Su arquitectura se compone de tres elementos:
un \textit{broker} que transmite los mensajes entre el productor y los
\textit{workers}, un \textit{result backend} opcional para almacenar los
resultados de las tareas, y los propios \textit{workers}. Como
\textit{broker} y \textit{result backend} el SCDEE emplea Redis
\cite{redis}, una base de datos en memoria con soporte para estructuras
de datos eficientes y operaciones atómicas, que cubre adicionalmente el
rol de caché de aplicación y de almacén de la lista negra de tokens. La
elección de Redis frente a alternativas robustas como RabbitMQ obedece
principalmente a su menor complejidad operativa: Redis basta sobradamente
para el volumen previsto de tareas y consolida en un único servicio
varias funciones de infraestructura, simplificando el despliegue.
